% !TeX spellcheck = es_ES
\documentclass[letterpaper,11pt]{article}
\usepackage[utf8]{inputenc}
\usepackage[T1]{fontenc}
\usepackage[spanish]{babel}
\decimalpoint % separador decimal con punto, como en las tablas de la fuente original
\usepackage{lmodern,microtype}
\usepackage{amsmath,amssymb,amsfonts}
\usepackage{geometry}
\geometry{margin=2.54cm}
\usepackage{setspace}
\setstretch{1.15}
\usepackage{enumitem}
\usepackage{booktabs}
\usepackage{color}
\usepackage{hyperref}
\definecolor{ntr}{rgb}{0.8,0,0}
\hypersetup{colorlinks=true,linkcolor=ntr,citecolor=ntr,urlcolor=black}

% Portada (Luis Cabral): titulo \LARGE sans-serif alineado a la izquierda, sin numero de pagina
\makeatletter
    \renewcommand{\@maketitle}{%
    \newpage\parindent0in\null\vskip2em%
    {\LARGE\textbf{\textsf{\@title}}\par\vskip2em}%
    \@author\par\vskip2em%
    \@date%
    \par
    \thispagestyle{empty}\setcounter{page}{1}%
    }%
\makeatother

\DeclareMathOperator{\Var}{Var}
\DeclareMathOperator{\Cov}{Cov}
\DeclareMathOperator{\Corr}{Corr}
\DeclareMathOperator{\E}{\mathbb{E}}
\DeclareMathOperator{\N}{Normal}
\DeclareMathOperator{\LP}{\mathbb{L}}
\newcommand\independent{\protect\mathpalette{\protect\independenT}{\perp}}
\def\independenT#1#2{\mathrel{\rlap{$#1#2$}\mkern2mu{#1#2}}}


\title{Econometría I \\ Ayudantía 1}
\author{Magíster en Economía, Universidad Alberto Hurtado}
\date{}

\begin{document}
\maketitle

\noindent Temas: probabilidad conjunta y marginal, esperanza y varianza, covarianza y correlación,
esperanza y varianza condicional, independencia e independencia en media, proyector lineal
(modelo de regresión lineal) y multicolinealidad perfecta.

\medskip
\noindent Los ejercicios marcados con \textbf{(*)} son los que alcanzamos a resolver en la
ayudantía (80 minutos); los demás quedan como práctica adicional, con solución en la pauta.

\begin{enumerate}[itemsep=1.2em]

%% ---------------------------------------------------------------- 1
\item \textbf{(*)} Dada la distribución de probabilidad de la siguiente tabla:
    \begin{table}[!htbp]
        \centering
        \caption{Distribución conjunta de condiciones meteorológicas y tiempo de viaje}
        \begin{tabular}{lcc}
            \toprule
            & Lluvia ($x=0$) & Sin lluvia ($x=1$) \\
            \midrule
            Viaje largo ($y=0$) & 0.15 & 0.07 \\
            Viaje corto ($y=1$) & 0.15 & 0.63 \\
            \bottomrule
        \end{tabular}
    \end{table}

    Calcule:
    \begin{enumerate}
        \item $\E(y)$ y $\E(x)$.
        \item $\Var(y)$ y $\Var(x)$.
        \item $\Cov(x,y)$ y la correlación entre $x$ e $y$.
    \end{enumerate}

%% ---------------------------------------------------------------- 2
\item \textbf{(*)} Usando las variables aleatorias $x$ e $y$ de la tabla anterior, considere dos nuevas
    variables aleatorias $w = 3 + 6\,x$ y $v = 20 - 7\,y$. Calcule:
    \begin{enumerate}
        \item $\E(w)$ y $\E(v)$.
        \item $\Var(w)$ y $\Var(v)$.
        \item $\Cov(w,v)$ y la correlación entre $w$ y $v$. Compare con el resultado del
              ejercicio 1(c) y explique qué cambió y qué no.
    \end{enumerate}

%% ---------------------------------------------------------------- 3
\item En septiembre, la temperatura máxima diaria en Seattle tiene una media de 70$^\circ$F y
    una desviación estándar de 7$^\circ$F. ¿Cuál es la media, la desviación estándar y la
    varianza en $^\circ$C? \\
    Recuerde que $^\circ\text{C} = (^\circ\text{F} - 32)\times 5/9$.

%% ---------------------------------------------------------------- 4
\item \emph{Esperanza y varianza condicional.} $x$ puede tomar dos valores $0,1$ e $y$ puede
    tomar tres valores $-1,0,1$, con la siguiente distribución conjunta:
    \begin{table}[!htbp]
        \centering
        \caption{Función de probabilidad conjunta}
        \begin{tabular}{rrcc}
            \toprule
            & & \multicolumn{2}{c}{$x$} \\
            \cmidrule(lr){3-4}
            & & 0 & 1 \\
            \midrule
                & $-1$ & 0.10 & 0.20 \\
            $y$ & $0$  & 0.30 & 0.10 \\
                & $1$  & 0.10 & 0.20 \\
            \bottomrule
        \end{tabular}
    \end{table}
    \begin{enumerate}
        \item Calcule las probabilidades marginales de $x$ e $y$, y $\E(y)$.
        \item Calcule las probabilidades condicionales de $y$ dado $x$ y $\E(y|x)$.
        \item Calcule $\Var(y|x=0)$ y $\Var(y|x=1)$.
        \item Verifique la descomposición de varianza
              $\Var(y) = \Var(\E(y|x)) + \E(\Var(y|x))$.
        \item ¿Son $x$ e $y$ independientes en media? ¿Son independientes? Concluya.
    \end{enumerate}

%% ---------------------------------------------------------------- 5
\item \textbf{(*)} \emph{Covarianza e independencia en media.} Suponga que las variables aleatorias $x$ y
    $z$ se distribuyen $\N(0,1)$ y son independientes, y suponga que $y = x^2 + z$. Encuentre
    la esperanza condicional de $y$ dado $x$ ($\E(y|x)$) y la covarianza entre $x$ e $y$
    ($\Cov(x,y)$). ¿Qué concluye? \\
    Pista: $\E[(x-\E x)^3] = 0$ para distribuciones simétricas como la normal.

%% ---------------------------------------------------------------- 6
\item \textbf{(*)} \emph{Independencia, independencia en media y correlación.} Suponga una variable
    aleatoria $u$ con media cero, $\E(u)=0$. Indique si la afirmación es verdadera o falsa y
    justifique (con una demostración o con un contraejemplo).
    \begin{enumerate}
        \item Si $\E(u|x)=0$ entonces $\E(x^2u)=0$.
        \item Si $\E(x\,u)=0$ entonces $\E(x^2u)=0$.
        \item Si $\E(u|x)=0$ entonces $u$ es independiente de $x$.
        \item Si $\E(x\,u)=0$ entonces $\E(u|x)=0$.
    \end{enumerate}

%% ---------------------------------------------------------------- 7
\item \emph{Aplicación.} Dada la siguiente distribución conjunta de situación laboral y
    estudios universitarios, para personas mayores de 25 años que trabajan o buscan trabajo,
    EEUU 2008.
    \begin{table}[!htbp]
        \centering
        \caption{Situación laboral y estudios universitarios}
        \begin{tabular}{rrccc}
            \toprule
            & & Desempleado & Empleado & \\
            & & ($y=0$) & ($y=1$) & Marginal de $x$ \\
            \midrule
            No titulado universitario & ($x=0$) & 0.037 & 0.622 & 0.659 \\
            Titulado universitario    & ($x=1$) & 0.009 & 0.332 & 0.341 \\
            \midrule
            Marginal de $y$ & & 0.046 & 0.954 & 1.000 \\
            \bottomrule
        \end{tabular}
    \end{table}
    \begin{enumerate}
        \item Calcule $\E(y)$.
        \item La tasa de desempleo es la proporción de la fuerza laboral que se encuentra
              desempleada. Demuestre que la tasa de desempleo está dada por $1-\E(y)$.
        \item Calcule $\E(y|x=1)$ y $\E(y|x=0)$.
        \item Calcule $\Var(y|x=1)$ y $\Var(y|x=0)$.
        \item Calcule la tasa de desempleo para (i) titulados y (ii) no titulados universitarios.
        \item Una persona de esta población seleccionada al azar dice estar desempleada.
              ¿Cuál es la probabilidad de que sea titulada universitaria? ¿Y no titulada?
        \item ¿Son independientes el nivel educativo y la situación laboral?
    \end{enumerate}

%% ---------------------------------------------------------------- 8
\item \textbf{(*)} \emph{Multicolinealidad perfecta.} Suponga que el vector de regresores,
    usado para predecir un escalar $y$, es
    \begin{align*}
    x = \begin{pmatrix} 1 \\ x_1 \\ x_2 \end{pmatrix},
    \end{align*}
    y que $x_2 = \alpha_1 + \alpha_2\, x_1$, es decir, $x_2$ es una función lineal exacta de $x_1$.
    \begin{enumerate}
        \item Demuestre que $\E(xx')$ no es invertible. ¿Cuál es la causa de la no invertibilidad?
        \item Sea $\beta$ un vector que minimiza $\E[(y-x'b)^2]$ y sea $c$ el vector no
              nulo de la parte (a). Demuestre que $\beta+t\,c$ predice exactamente el mismo $y$
              que $\beta$, para cualquier $t\in\mathbb{R}$, y que por lo tanto también minimiza
              $\E[(y-x'b)^2]$. ¿Qué consecuencias tiene esto para la identificación de $\beta$ en
              el modelo de regresión lineal $\LP(y|x)=x'\beta$?
    \end{enumerate}

%% ---------------------------------------------------------------- 9
\item \textbf{(*)} \emph{Proyección lineal con solo un intercepto.} Considere el modelo
    $y=\beta_0+\varepsilon$, donde $\beta_0$ es el proyector lineal, es decir
    $\beta_0=\arg\min_b \E[(y-b)^2]$. Demuestre que $\beta_0=\E(y)$.

\end{enumerate}
\end{document}
